\documentclass[11pt,a4paper]{article}

\usepackage[margin=2.4cm]{geometry}
\usepackage[T1]{fontenc}
\usepackage{lmodern}
\usepackage{xcolor}
\usepackage{graphicx}
\usepackage{booktabs}
\usepackage{titlesec}
\usepackage{enumitem}
\usepackage{fancyhdr}
\usepackage{listings}
\usepackage{longtable}
\usepackage[hidelinks]{hyperref}

\definecolor{aimzblue}{HTML}{0B5FFF}
\definecolor{codebg}{HTML}{F5F6F8}
\definecolor{codeframe}{HTML}{D8DCE3}
\definecolor{warnbg}{HTML}{FDECEA}
\definecolor{warnframe}{HTML}{C0392B}

\hypersetup{
  colorlinks=true,
  linkcolor=aimzblue,
  urlcolor=aimzblue,
  pdftitle={AIMZ Egypt --- Cloudflare to AWS Migration Runbook},
  pdfauthor={AIMZ Platform}
}

\lstset{
  backgroundcolor=\color{codebg},
  basicstyle=\ttfamily\small,
  frame=single,
  rulecolor=\color{codeframe},
  breaklines=true,
  breakatwhitespace=true,
  columns=fullflexible,
  keepspaces=true,
  showstringspaces=false,
  xleftmargin=6pt,
  xrightmargin=6pt,
  aboveskip=8pt,
  belowskip=8pt,
  literate={~}{{\raisebox{0.5ex}{\texttildelow}}}1
}

\titleformat{\section}{\Large\bfseries\color{aimzblue}}{\thesection}{0.6em}{}
\titleformat{\subsection}{\large\bfseries}{\thesubsection}{0.6em}{}

\pagestyle{fancy}
\fancyhf{}
\lhead{\small AIMZ Egypt}
\rhead{\small Cloudflare $\rightarrow$ AWS Migration Runbook}
\cfoot{\thepage}
\renewcommand{\headrulewidth}{0.4pt}

% Warning box
\newcommand{\warnbox}[1]{%
  \begingroup
  \setlength{\fboxsep}{8pt}%
  \noindent\colorbox{warnbg}{\parbox{\dimexpr\linewidth-2\fboxsep}{%
  \color{warnframe}\bfseries #1}}%
  \endgroup\par\vspace{6pt}
}

\begin{document}

\begin{titlepage}
\centering
\vspace*{3cm}
{\Huge\bfseries\color{aimzblue} AIMZ Egypt\par}
\vspace{0.6cm}
{\LARGE Cloudflare $\rightarrow$ AWS Migration Runbook\par}
\vspace{0.4cm}
{\large How to launch the platform and push updates to AWS\par}
\vspace{2cm}
{\large Infrastructure: AWS CDK (TypeScript)\par}
{\large Compute: EC2 + ALB \quad Data: RDS PostgreSQL + ElastiCache Redis\par}
{\large Storage/CDN: S3 + CloudFront\par}
\vfill
{\large Target scale: \textbf{$\sim$500 concurrent users}\par}
\vspace{0.5cm}
{\small Generated for the \texttt{aimz} monorepo --- \texttt{infra/aws-cdk}\par}
\end{titlepage}

\tableofcontents
\newpage

% ---------------------------------------------------------------------------
\section{Read this first: security}

\warnbox{An AWS access key was shared in plaintext during preparation of this
migration. Treat it as compromised. Before anything else: sign in to the AWS
console, open IAM $\rightarrow$ Users $\rightarrow$ Security credentials,
deactivate and delete that access key, then create a fresh one. Never paste
long-lived keys into chats, tickets, or source files.}

This project never stores AWS credentials in code. The CDK and the helper
scripts authenticate through your local AWS CLI profile
(\texttt{aws configure} or SSO). Application secrets (database password, JWT
secret, admin password, SMTP) live only in \textbf{AWS Secrets Manager} and are
pulled by the API instances at boot.

% ---------------------------------------------------------------------------
\section{What this migration does}

The live API today runs on Cloudflare Workers (Hono) backed by D1 (SQLite) and
R2, with the web app on Cloudflare Pages. This migration moves the platform onto
a conventional, horizontally-scalable AWS stack driven entirely by the FastAPI
backend in \texttt{backend/}, which already speaks PostgreSQL and S3.

\begin{center}
\begin{tabular}{@{}lll@{}}
\toprule
\textbf{Concern} & \textbf{Cloudflare (before)} & \textbf{AWS (after)} \\
\midrule
API compute      & Workers (Hono)        & EC2 Auto Scaling Group + ALB \\
Database         & D1 (SQLite)           & RDS PostgreSQL \\
Cache            & --- (none)            & ElastiCache Redis \\
Object storage   & R2                    & S3 (private, presigned) \\
Web hosting      & Pages                 & S3 + CloudFront (OAC) \\
Secrets          & Worker vars           & Secrets Manager \\
Image registry   & ---                   & ECR (via CDK assets) \\
\bottomrule
\end{tabular}
\end{center}

\subsection{Architecture at a glance}

\begin{lstlisting}[language=]
                    Internet
                       |
        +--------------+---------------+
        |                              |
   CloudFront (web)               ALB  (api, :443/:80)
        |                              |
   S3 web bucket             EC2 ASG (FastAPI :8000, 2..4 x t3.medium)
   (private, OAC)                 |            |
                                  |            +--> S3 media bucket (presigned)
                                  |
                      +-----------+-----------+
                      |                       |
               RDS PostgreSQL         ElastiCache Redis
               (private isolated)     (private isolated)
\end{lstlisting}

All data services live in private, isolated subnets. Only the ALB and
CloudFront are internet-facing. The API tier reaches the internet for image
pulls and AWS APIs through a single NAT gateway.

% ---------------------------------------------------------------------------
\section{Prerequisites}

\begin{itemize}[leftmargin=1.4em]
  \item An AWS account and an IAM identity with administrator (or equivalently
        broad) permissions for the initial bootstrap.
  \item Local tools: \texttt{aws} CLI v2, Node.js 20+, Docker (running),
        \texttt{jq}, and \texttt{git}.
  \item AWS CLI configured for the target account and region:
\begin{lstlisting}[language=bash]
aws configure          # or: aws configure sso
aws sts get-caller-identity   # confirm you are the right account
export AWS_REGION=eu-west-1    # pick your region once
\end{lstlisting}
  \item Node dependencies installed for the CDK app:
\begin{lstlisting}[language=bash]
cd infra/aws-cdk
npm install
\end{lstlisting}
\end{itemize}

% ---------------------------------------------------------------------------
\section{One-time bootstrap}

CDK needs a small bootstrap stack (an asset bucket, ECR repo, and roles) in each
account/region before the first deploy:

\begin{lstlisting}[language=bash]
cd infra/aws-cdk
npx cdk bootstrap aws://<ACCOUNT_ID>/<REGION>
\end{lstlisting}

Run this once per account/region. It is safe to re-run.

% ---------------------------------------------------------------------------
\section{Configuration knobs}

Everything tunable is passed as CDK \emph{context} (\texttt{-c key=value}) or an
environment variable, with defaults sized for $\sim$500 concurrent users.

\begin{center}
\begin{longtable}{@{}lll@{}}
\toprule
\textbf{Context key} & \textbf{Default} & \textbf{Purpose} \\
\midrule
\endhead
\texttt{env}               & \texttt{production} & Logical environment / stack suffix \\
\texttt{apiInstanceType}   & \texttt{t3.medium}  & EC2 class for API instances \\
\texttt{apiMinCapacity}    & \texttt{2}          & Min instances (HA across 2 AZs) \\
\texttt{apiMaxCapacity}    & \texttt{4}          & Max instances (CPU target 60\%) \\
\texttt{dbInstanceClass}   & \texttt{t3.medium}  & RDS instance class \\
\texttt{dbMultiAz}         & \texttt{false}      & RDS Multi-AZ standby (set true for prod HA) \\
\texttt{redisNodeType}     & \texttt{cache.t4g.small} & ElastiCache node class \\
\texttt{webDomainName}     & (none)              & Custom domain for CloudFront \\
\texttt{webCertificateArn} & (none)              & ACM cert (us-east-1) for the web domain \\
\texttt{apiCertificateArn} & (none)              & ACM cert (stack region) for ALB HTTPS \\
\bottomrule
\end{longtable}
\end{center}

\subsection{Sizing note for 500 concurrent users}

The defaults give two \texttt{t3.medium} API nodes (2 vCPU / 4\,GiB each) behind
an ALB, auto-scaling to four at 60\% CPU. A scores/roster workload is bursty and
read-heavy, so a \texttt{db.t3.medium} with the Redis cache in front absorbs
500 concurrent users comfortably. For sustained peak traffic, raise
\texttt{apiMaxCapacity}, move the database to \texttt{db.m6g.large}, and set
\texttt{dbMultiAz=true}.

% ---------------------------------------------------------------------------
\section{Launch (first deploy)}

From \texttt{infra/aws-cdk}:

\begin{lstlisting}[language=bash]
# 1. Review what will be created (optional but recommended)
npx cdk diff -c env=production

# 2. Create everything. CDK builds the backend Docker image, pushes it to ECR,
#    and stands up VPC, RDS, Redis, S3, CloudFront, ALB and the ASG.
npx cdk deploy -c env=production --require-approval never
\end{lstlisting}

The first deploy takes roughly 20--30 minutes (RDS and CloudFront dominate). At
the end, CDK prints the stack outputs; the ones you need:

\begin{center}
\begin{tabular}{@{}ll@{}}
\toprule
\textbf{Output} & \textbf{Use} \\
\midrule
\texttt{ApiUrl}   & Public API base URL (goes into the mobile build) \\
\texttt{WebUrl}   & CloudFront URL of the web app \\
\texttt{AsgName}  & Used by the migration/redeploy scripts \\
\texttt{DbEndpoint}, \texttt{RedisEndpoint} & Data endpoints (informational) \\
\bottomrule
\end{tabular}
\end{center}

\subsection{Post-deploy steps (run once)}

\begin{enumerate}[leftmargin=1.6em]
  \item \textbf{Set the admin password.} Hosted environments refuse the
        placeholder value:
\begin{lstlisting}[language=bash]
./scripts/set-admin-password.sh production 'a-strong-admin-password' admin@yourdomain.com
./scripts/redeploy-api.sh production   # recycle instances to pick it up
\end{lstlisting}
  \item \textbf{Run database migrations + seed.} This applies Alembic migrations
        and creates the admin account and initial invite. It runs inside one
        instance's container over SSM (no SSH), so it never races the ASG:
\begin{lstlisting}[language=bash]
./scripts/run-migrations.sh production
\end{lstlisting}
  \item \textbf{Publish the web app.} Builds the Expo web export against
        \texttt{ApiUrl} and pushes it to the CloudFront bucket:
\begin{lstlisting}[language=bash]
./scripts/deploy-web.sh production
\end{lstlisting}
  \item \textbf{Verify.}
\begin{lstlisting}[language=bash]
curl -s "$(aws cloudformation describe-stacks --stack-name Aimz-production \
  --query "Stacks[0].Outputs[?OutputKey=='ApiUrl'].OutputValue" --output text)/api/v1/health"
# -> {"status":"ok","service":"aimz-api",...}
\end{lstlisting}
\end{enumerate}

% ---------------------------------------------------------------------------
\section{Pushing updates}

\subsection{Backend / API changes}

Any change under \texttt{backend/} ships as a new container image:

\begin{lstlisting}[language=bash]
cd infra/aws-cdk
npx cdk deploy -c env=production     # rebuilds + pushes the image to ECR
./scripts/redeploy-api.sh production # rolling instance refresh onto the new image
# If the change added migrations:
./scripts/run-migrations.sh production
\end{lstlisting}

\subsection{Web / mobile-web changes}

\begin{lstlisting}[language=bash]
cd infra/aws-cdk
./scripts/deploy-web.sh production    # rebuild export, sync to S3, invalidate CDN
\end{lstlisting}

\subsection{Infrastructure changes}

Edit \texttt{lib/aimz-stack.ts}, then:

\begin{lstlisting}[language=bash]
npx cdk diff -c env=production        # preview
npx cdk deploy -c env=production      # apply
\end{lstlisting}

% ---------------------------------------------------------------------------
\section{Operations}

\begin{description}[leftmargin=0em, style=nextline]
  \item[Shell into an instance] \texttt{aws ssm start-session --target <instance-id>}
        (no SSH keys, no open ports). Then \texttt{docker logs -f aimz-api}.
  \item[API logs] The container logs to the instance; for aggregation, ship
        \texttt{docker logs} to CloudWatch (future enhancement) or use the
        Session Manager shell above.
  \item[Scaling] The ASG target-tracks CPU at 60\%. Change
        \texttt{apiMin/MaxCapacity} and redeploy, or edit the ASG in the console
        for an immediate bump.
  \item[Database backups] RDS automated backups retain 7 days. Production uses a
        final snapshot on stack deletion and deletion protection.
  \item[Secrets rotation] Update the value in Secrets Manager (or via
        \texttt{set-admin-password.sh}) and run \texttt{redeploy-api.sh} so
        instances re-read it at boot.
\end{description}

% ---------------------------------------------------------------------------
\section{Cost sketch (order of magnitude)}

Rough monthly, on-demand, single region --- confirm against the AWS pricing
calculator for your region:

\begin{center}
\begin{tabular}{@{}lr@{}}
\toprule
\textbf{Component} & \textbf{Est. / month (USD)} \\
\midrule
2$\times$ t3.medium EC2 (API)        & \textasciitilde{}60 \\
ALB                                  & \textasciitilde{}20 \\
RDS db.t3.medium (single-AZ) + 50GB  & \textasciitilde{}70 \\
ElastiCache cache.t4g.small          & \textasciitilde{}25 \\
NAT gateway                          & \textasciitilde{}35 \\
S3 + CloudFront (light)              & \textasciitilde{}10 \\
\midrule
\textbf{Total}                       & \textbf{\textasciitilde{}220} \\
\bottomrule
\end{tabular}
\end{center}

Multi-AZ RDS and a Redis replica roughly double their respective lines and are
recommended before you rely on this in production.

% ---------------------------------------------------------------------------
\section{Teardown}

\warnbox{Destroys data. In production the database takes a final snapshot and
buckets are retained; elsewhere everything is removed.}

\begin{lstlisting}[language=bash]
cd infra/aws-cdk
npx cdk destroy -c env=production
\end{lstlisting}

% ---------------------------------------------------------------------------
\section{Remaining work --- stood down}

\textbf{This migration is not being pursued.} Cloudflare is the platform:
Pages serves the web app, a Worker serves the API, and D1 holds the data.
This runbook covers \textbf{Stage 1} only --- the AWS infrastructure and
deployment tooling --- and Stage 1 is where it stopped. Nothing has been
deleted, so the option survives, but neither stage below is being worked on.

\begin{enumerate}[leftmargin=1.6em]
  \item \textbf{Stage 2 --- FastAPI feature parity. Never completed.} An
        earlier revision of this page put the gap at ``roughly 15 migrations and
        eight route groups (stats, knockout, training, announcements, roster,
        calendar, audit)''. That estimate was wrong in both directions and
        should not be relied on: \texttt{backend/} already had every one of
        those groups, while the families it actually lacks were missing from
        the list entirely. As of this writing \texttt{backend/} serves 14 of
        the 26 endpoint families the mobile app calls, and is 18 Alembic
        migrations behind the Worker's 46 D1 migrations. The twelve it does
        not serve are fee plans, fee charges, fee payments, fee invoices,
        invoices, kit orders, match reports, player reports, attendance
        requests, branches, shared reports, and training metrics.

        Do not copy a number out of this paragraph. Run \texttt{npm run parity}
        in \texttt{cloudflare-api/}; it prints the current gap, and it exists
        because every hand-typed parity figure in this repository has gone
        stale, including the one above.
  \item \textbf{Stage 3 --- Cutover. Not safe to run.} See
        \texttt{docs/aws-migration/stage3-data-cutover.md}, which is marked
        stood down. It once claimed Stage 2 had achieved parity; it had not,
        and following it would have moved live academy data onto an API that
        cannot serve half the product.
\end{enumerate}

\vfill
\begin{center}\small
\textit{Infrastructure as code: \texttt{infra/aws-cdk/lib/aimz-stack.ts} ---
one \texttt{cdk deploy} stands up the whole platform.}
\end{center}

\end{document}
